\section{Introduction}
\label{sec:intro}

Consumer Internet-of-Things devices are numerous, rarely patched, and
frequently exposed on home networks with default credentials. The Mirai botnet
showed the consequence at scale: a large population of cameras, routers and
recorders was conscripted into denial-of-service attacks largely by guessing
factory passwords~\cite{antonakakis2017mirai}. A household has no security
operations team to watch its traffic. Whatever defends it has to run
unattended, on modest hardware, and without shipping private traffic to a
third party.

This paper describes \textsc{Sentinel}, an open-source appliance that attempts
exactly that, and reports what a careful evaluation of it revealed.
\textsc{Sentinel} captures traffic on the home segment, scores it with an
autoencoder over Kitsune's streaming features~\cite{mirsky2018kitsune}, asks a
locally hosted large language model (LLM) to label anomalous events with MITRE
ATT\&CK tactics and techniques~\cite{mitreattack}, and turns the label into an
nftables~\cite{nftables} block through a deterministic, auditable response
policy. Honeypots~\cite{cowrie,opencanary} feed attacker addresses into the
same enforcement path. The language model runs on the appliance itself, so no
packet payload or traffic content leaves the home.

Building such a pipeline is now easy; knowing whether it works is not. Anomaly
detectors for IoT traffic are routinely reported at near-perfect accuracy on
public benchmarks, yet aggregate scores on lab datasets say little about how a
system behaves when deployed~\cite{sommer2010closedworld,arp2022dosdonts}. We
therefore treated evaluation as the main object of study, and it overturned
several things we believed about our own system:

\begin{itemize}
\item \emph{A per-class blind spot hidden by the aggregate.} On
  N-BaIoT~\cite{meidan2018nbaiot}, the deployed detector catches the Mirai
  attacks almost perfectly, yet misses nearly all BASHLITE TCP and UDP floods,
  which make up \NblindShare{} of attack windows. An aggregate
  detection rate conceals this entirely.
\item \emph{False positives grow on unseen devices.} Holding out entire devices
  from training, as a new gadget joining a home network would be, raises the
  false-positive rate from \NlstmFPRin{} to \NlstmFPRout{} for our
  reproducible model.
\item \emph{The live feature extractor did not match the training features.}
  \textsc{Sentinel}'s own re-implementation of the Kitsune extractor computed
  40 of the 115 features differently from the data the model was trained on,
  so offline accuracy did not describe the running system. We traced the
  discrepancy, fixed it by adopting the reference implementation, and verified
  parity against invariants measured on the training data itself.
\item \emph{The LLM's labels were almost never right.} In a pre-registered
  test, the local model named a correct ATT\&CK technique for
  \NllmAsbTech{} of attack windows, labelling \NllmCtwoShare{} of all
  events as Command and Control; describing the traffic to it did not help.
  Because the response policy acts on the model's severity, which it set high
  for most events, the pipeline would still have blocked \NllmAsbBlock{} of
  attacks---for reasons unrelated to the label.
\end{itemize}

\noindent Our contributions are:
\begin{enumerate}
\item The design and open-source implementation of an end-to-end,
  on-premise detection-to-response pipeline for home IoT networks that keeps
  inference, including the LLM, on the appliance (\S\ref{sec:design}).
\item A leakage-aware evaluation protocol for N-BaIoT with a
  \emph{held-out-device} split, applied to the deployed model, a reproducible
  retraining over five seeds, and two baselines (\S\ref{sec:method},
  \S\ref{sec:res-detect}).
\item A pre-registered evaluation of LLM ATT\&CK labelling on the events the
  pipeline actually produces, with and without a traffic summary, showing that
  a small local model does not yet perform this task (\S\ref{sec:res-llm}).
\item A per-stage latency breakdown on a four-core aarch64 host, locating the
  bottleneck of the detection-to-response path (\S\ref{sec:res-latency}).
\item Practical lessons for building and evaluating such systems: recovering
  the provenance of an undocumented model from its scaler, and checking that
  live features match training features (\S\ref{sec:discussion}).
\end{enumerate}

All code, the evaluation harness, the pre-registration and the result files
are released with the system.\footnote{Repository URL withheld for review.}
